Preface

Buddhism should not be considered a teaching for monks only. In the course of His ministry for forty-five years, the Buddha was indefatigable. He traveled on foot with a company of monks all over Northern India, from Vesali in the east to Kuru (Delhi) in the west, preaching the Dhamma for the benefit of mankind. Those who wished to renounce the household life to devote their entire life to the study and practice of the Dhamma entered the Sangha and became monks or nuns. Those who could not renounce the household life practised as lay followers and many attained to various stages the sainthood as recorded in the Pali scriptures.

The main problem for the lay Buddhist is how to combine personal progress in worldly matters with moral principles. During his time and even nowadays, the majority of people who choose Buddhism as their religion would want to know what are the qualities that one should develop to become a good Buddhist! So the compassionate Buddha preached the Candala Sutta in the Anguttara Nikaya V, 175, to the Bikkhus (monks) so that they can propagate his instructions to the people even after his demise or Parinibbhana.

Among the discourses for the laity, the Candala Sutta or Discourse on Outcaste is perhaps the simplest and most suitable discourse for the newcomer to understand and follow the Buddha’s Teachings. It defines five qualities that distinguish a Lotus-like lay follower from an outcaste in the Buddha’s Dispensation (Sasana), namely: Faith in the Buddha, Dhamma, and Sangha; Morality; Absence of superstitious belief; Trust in the Law of Kamma; Seeking with (the Sangha) for a gift-worthy person and offering here first.

In the course of teaching this sutta at several Buddhist societies in the Klang Valley, the author found that it was well received by an appreciative audience who clearly understood the essentials taught in this sutta. They could see that there is no mystery or mystique in the Buddha’s Teaching, which is based on the Four Noble Truths and Dependent Origination. In order to help the lay Buddhist to understand through which causes and conditions suffering arises, a chapter on the Doctrine of Dependent Origination has been added.

During his ministry, Lord Buddha taught two kinds of Right View or Understanding, namely: that pertaining to the worldly and that pertaining to Enlightenment. Of the two views, the former is meant the right road to the round of existences (Samsara) while the latter is the right road out of the rounds of existence (Nibbana). The worldly or mundane Right View is the Understanding of the Ten Subjects and Law of Kamma. To understand the Four Noble Truths, which is the Teaching of the Buddha, one must meditate upon one’s mentality-materiality (mind-body) by Satipatthana Vipassana meditation to penetrate their true nature, viz. impermanence, suffering and non-self. This insight knowledge will temporarily free one from wrong view of a personality or ‘sakkaya-ditthi’ in Pali.

In Anguttara Nikaya II, 32, the Buddha said that no matter how many material gifts we give to our parents, we can never repay our debt of gratitude to them. All the material things we provide them now can only benefit them in this life. But if we can establish them in faith, morality, generosity and wisdom, it would not only benefit them now but also in future lives. With this Gift of Dhamma, one can really repay one’s debt of gratitude to parents!

I take this opportunity to thank Sis Wooi Kheng Choo for checking the texts and giving useful suggestions.

Bro. Chan Khoon San, Klang, 1 July 2012

